% Source: https://artofproblemsolving.com/community/c3h1574785_1011_nice
% Stable identity: aops-topic-1574785-q01
Let the parts, in order, have lengths $a,b,c,d$ with $a+b+c+d=36$, and let $M_1,\dots,M_4$ be their midpoints.

\begin{asy}
usepackage("amsmath");
size(8cm);
defaultpen(fontsize(10pt));
real[] p={0,5,13,29,36};
draw((0,0)--(36,0),linewidth(0.9));
for(int i=0;i<5;++i) draw((p[i],-0.6)--(p[i],0.6));
pair[] M; for(int i=0;i<4;++i) M.push(((p[i]+p[i+1])/2,0));
for(int i=0;i<4;++i) dot(M[i],red);
label("$a$",((p[0]+p[1])/2,0.6),N);label("$b$",((p[1]+p[2])/2,0.6),N);label("$c$",((p[2]+p[3])/2,0.6),N);label("$d$",((p[3]+p[4])/2,0.6),N);
label("$M_1$",M[0],S,red);label("$M_2$",M[1],S,red);label("$M_3$",M[2],S,red);label("$M_4$",M[3],S,red);
draw((M[0].x,-3)--(M[3].x,-3),blue,Arrows(TeXHead));label("$30$",((M[0].x+M[3].x)/2,-3),S,blue);
draw((M[1].x,-5.2)--(M[2].x,-5.2),deepgreen,Arrows(TeXHead));label("$?$",((M[1].x+M[2].x)/2,-5.2),S,deepgreen);
draw((M[0].x,-0.9)--(M[0].x,-3),gray+dotted);draw((M[3].x,-0.9)--(M[3].x,-3),gray+dotted);
draw((M[1].x,-0.9)--(M[1].x,-5.2),gray+dotted);draw((M[2].x,-0.9)--(M[2].x,-5.2),gray+dotted);
label("\mbox{}",(min(currentpicture,user=true)+max(currentpicture,user=true))/2,black); // keeps the text after the figure black when LaTeX typesets the labels inline
\end{asy}

(The figure shows one possible division, $5,8,16,7$, not to scale for the general case.)

From $M_1$ to $M_4$ we cover half of the first part, all of the second and third parts, and half of the fourth:
\[
M_1M_4=\dfrac {a}{2}+b+c+\dfrac {d}{2}=36-\dfrac{a+d}{2}=30,\qquad\text{so}\qquad a+d=12,\quad b+c=24 .
\]
From $M_2$ to $M_3$ we cover half of the second part and half of the third:
\[
M_2M_3=\dfrac {b}{2}+\dfrac {c}{2}=\dfrac{24}{2}=\boxed{12\text{ cm}}.
\]
The answer does not depend on how the segment is divided (for example $5,8,16,7$ or $3,10,14,9$).
